\begin{pastexam}[label={px:df3-104-ds-02}]{地方三等資訊處理 104 年 第二題}
（申論；20 分）

請使用 C 或 Java 語言寫一副程式 void merge(int [] A, int [] B, int [] C, int n)，此副程式將對
兩個長度為 n 且已依小到大排序的整數陣列 A 與 B，合併至長度為 2n 且依小到大排序的整數陣列
C，此副程式的時間複雜度需為 $\Theta$(n)。（20 分）

（注意：請加註解說明程式碼作法）
\end{pastexam}

\begin{pastsoln}
題目的參數寫法 \texttt{int [] A} 是 Java 的語法；用 C 寫時改成 \texttt{int A[]}，意思相同。

\textbf{作法。} 兩個陣列都已排好，最小的數一定是 $A[0]$ 或 $B[0]$ 之一。 用 $i$、$j$ 分別指著
$A$、$B$ 目前還沒放進 $C$ 的最前面那一格，每次比較 $A[i]$ 和 $B[j]$，把小的放進 $C[k]$，
該指標往後一格。 其中一個陣列用完後，另一個剩下的部分本來就排好、而且都不比 $C$ 裡已放的數
小，整段照抄到 $C$ 的後面（資結書第~14.4~節的合併）。
\begin{minted}[fontsize=\small, style=perldoc]{c}
/* 把已由小到大排好的 A[0..n-1]、B[0..n-1] 合併到 C[0..2n-1] */
void merge(int A[], int B[], int C[], int n)
{
    int i = 0;      /* A 中下一個還沒放進 C 的位置 */
    int j = 0;      /* B 中下一個還沒放進 C 的位置 */
    int k = 0;      /* C 中下一個要填的位置 */

    /* 兩邊都還有數時，取兩邊最前面較小的那一個 */
    while (i < n && j < n) {
        if (A[i] <= B[j])       /* 相等時取 A 的，合併是穩定的 */
            C[k++] = A[i++];
        else
            C[k++] = B[j++];
    }
    /* 其中一邊用完了，另一邊剩下的已排好，依序抄過去 */
    while (i < n)
        C[k++] = A[i++];
    while (j < n)
        C[k++] = B[j++];
}
\end{minted}

\textbf{正確性}：迴圈不變量是「$C[0..k-1]$ 是 $A[0..i-1]$ 和 $B[0..j-1]$ 合起來由小到大排好的
結果，且其中每個數都 $\le A[i]$ 和 $B[j]$（存在的話）」。 一開始 $k = 0$，成立。 第一個迴圈每次
放進的是 $A[i]$、$B[j]$ 中較小的，它 $\ge C[k-1]$（由不變量），又 $\le$ 兩邊剩下的所有數（兩邊
都排好了），所以放進去之後不變量仍然成立；後兩個迴圈只剩一邊，同理。 結束時 $i = j = n$、
$k = 2n$，$C[0..2n-1]$ 是全部 $2n$ 個數排好的結果。

\textbf{追蹤}：$A = \langle 1, 4, 6, 9 \rangle$、$B = \langle 2, 3, 7, 10 \rangle$，$n = 4$。
\begin{center}
\small
\begin{tabular}{cccl}
\toprule
比較 & 放進 $C$ & $(i, j)$ 之後 & $C$ \\
\midrule
$1 \le 2$ & 1 & (1, 0) & 1 \\
$4 > 2$ & 2 & (1, 1) & 1, 2 \\
$4 > 3$ & 3 & (1, 2) & 1, 2, 3 \\
$4 \le 7$ & 4 & (2, 2) & 1, 2, 3, 4 \\
$6 \le 7$ & 6 & (3, 2) & 1, 2, 3, 4, 6 \\
$9 > 7$ & 7 & (3, 3) & 1, 2, 3, 4, 6, 7 \\
$9 \le 10$ & 9 & (4, 3) & 1, 2, 3, 4, 6, 7, 9 \\
（$A$ 用完） & 10 & (4, 4) & 1, 2, 3, 4, 6, 7, 9, 10 \\
\bottomrule
\end{tabular}
\end{center}

\textbf{時間複雜度}：每做一次迴圈本體，$k$ 就加一，$k$ 從 $0$ 到 $2n$，三個迴圈合計恰好執行
$2n$ 次，每次常數時間；比較次數最多 $2n - 1$ 次、最少 $n$ 次（一邊全部比另一邊小）。 所以時間
是 $\Theta(2n) = \Theta(n)$，最好、最壞都一樣，符合題目要求。 額外空間只有三個註標，$O(1)$
（$C$ 是呼叫端給的）。
\end{pastsoln}
